Interrupt sources are level-based and derived directly from kernel status
or FIFO fill levels.

Table~\ref{tab:asQspiPer06c} shows the service requests of the QSPI.
\begin{table}[H]
\caption{QSPI Service Requests (updated)}
\label{tab:asQspiPer06c}
\centering
\begin{tabularx}{\textwidth}{|l |l |l |X|}
  \hline
  Bit & Service Request & IRQ & Explanation \\
  \hline
  \hline
  0 & DONE      & DONE      & Transfer completed (FSM reached DONE state) \\
  \hline
  1 & RX\_HALF  & RX\_HALF  & RX FIFO $\geq$ half full \\
  \hline
  2 & RX\_EMPTY & RX\_EMPTY & RX FIFO empty \\
  \hline
  3 & TX\_HALF  & TX\_HALF  & TX FIFO $\leq$ half full (ready to be refilled) \\
  \hline
  4 & TX\_EMPTY & TX\_EMPTY & TX FIFO empty \\
  \hline
  5 & ERROR     & ERROR     & Error latched (see STATUS register) \\
  \hline
  6 & TIMEOUT   & TIMEOUT   & Timeout threshold reached \\
  \hline
  63:7 & -       & -         & Reserved \\
  \hline
\end{tabularx}
\end{table}

In QSPI: \[ \text{MIS} = \text{RIS} \;\&\; \text{IMSC} \]
and then the external interrupt output is: \[ \text{IRQ} = |\;\text{MIS} \]
